%! Author = denis
%! Date = 11/08/2026

\begin{frame}{Identidades lógicas: Lei da dupla negação}
    \begin{block}{Definição}
        Negar uma variável duas vezes consecutivas resulta no seu
        valor original.
    \end{block}

    \begin{columns}[T]

        % Coluna A = 0
        \begin{column}{0.5\textwidth}
            \begin{block}{Para $A = 0$}
                \[
                    A = 0
                \]
                \[
                    \overline{A} = 1
                \]
                \[
                    \overline{\overline{A}} = 0
                \]
                \[
                    \boxed{\overline{\overline{0}} = 0}
                \]
            \end{block}
        \end{column}

        % Coluna A = 1
        \begin{column}{0.5\textwidth}
            \begin{block}{Para $A = 1$}
                \[
                    A = 1
                \]
                \[
                    \overline{A} = 0
                \]
                \[
                    \overline{\overline{A}} = 1
                \]
                \[
                    \boxed{\overline{\overline{1}} = 1}
                \]
            \end{block}
        \end{column}

    \end{columns}

\end{frame}